\hypertarget{appendix-ceb4c4a9b5a1}{}\label{referencestart:appendix-ceb4c4a9b5a1}
\section{Fixed-point range and resolution}\label{appendix-ceb4c4a9b5a1-fixed-point-range-and-resolution}

\hypertarget{appendix-ceb4c4a9b5a1-fixed-point-range-and-resolution}{}

\textbf{Curriculum depth:} Developed · \textbf{First seen:} C006 · \textbf{Latest development:} C006

\begin{allmightroledefinition}

\subsection{Definition and mechanism}\label{appendix-ceb4c4a9b5a1-definition-and-mechanism}

\hypertarget{appendix-ceb4c4a9b5a1-definition-and-mechanism}{}

For signed N-bit raw storage and F fraction bits, step is 2\^{}-F and the positive endpoint is (2\textsuperscript{(N-1)-1)/2}F. Required step 0.001 and magnitude 1e6 in 32 bits allow F=10 or 11.

\end{allmightroledefinition}

\subsection{Prerequisites}\label{appendix-ceb4c4a9b5a1-prerequisites}

\hypertarget{appendix-ceb4c4a9b5a1-prerequisites}{}

\begin{itemize}
\tightlist
\item
  \hyperref[appendix-1d0dbe8924cc]{Fixed-point representation} --- supplies the representation or validity rule used above.
\end{itemize}

\subsection{Used by}\label{appendix-ceb4c4a9b5a1-used-by}

\hypertarget{appendix-ceb4c4a9b5a1-used-by}{}

Connect later dependent concepts here when their explanations use this mechanism.

\subsection{Recall and next depth}\label{appendix-ceb4c4a9b5a1-recall-and-next-depth}

\hypertarget{appendix-ceb4c4a9b5a1-recall-and-next-depth}{}

Reconstruct the example and identify the condition that would invalidate it.

\textbf{Pending:} Coordinate-origin and wider-carrier designs need their own range proofs.

\subsection{Source and navigation}\label{appendix-ceb4c4a9b5a1-source-and-navigation}

\hypertarget{appendix-ceb4c4a9b5a1-source-and-navigation}{}

\hyperref[01-foundations-i-ch001-choosing-what-information-to-keep]{Continuous teaching story} · {Source coverage ledger} · \href{../07-Vision-XR-and-Playable-Worlds-P01/main.pdf\#appendix-4ad47bdd269a}{Dependency Map (Vision, XR, and Playable Worlds, Part 1, Reference notes: Dependency Map)}

\subsection{Reconstruction depth}\label{appendix-ceb4c4a9b5a1-reconstruction-depth}

\hypertarget{appendix-ceb4c4a9b5a1-reconstruction-depth}{}

A fixed-point value stores an integer \(q\) together with an implied scale \(s\), so the represented real quantity is

\[
x=q\,s.
\]

For a binary fractional format with \(F\) fractional bits, \(s=2^{-F}\). Increasing \(F\) improves resolution but consumes integer range because storage width has not changed. For signed \(N\)-bit two's-complement storage, the approximate real range is

\[
[-2^{N-1}s,(2^{N-1}-1)s].
\]

This makes format choice an error-budget decision rather than a syntax decision. Position data may need large range; normalized coefficients may need fine local resolution. Multiplication temporarily doubles the number of fractional bits, so a production operation needs a widened intermediate or a deliberate rescale, round and saturate policy.

Tests should exercise representable endpoints, half-step rounding, overflow, negative values, multiplication rescaling and accumulation. A format can be perfectly deterministic and still be numerically unsuitable if its step size exceeds the physical or visual tolerance of the subsystem using it.
